\section*{Conventions and manuscript coordinates}
We distinguish a classifying topos of a discrete groupoid, a topos of continuous actions of a profinite groupoid, a pro-object of groupoids, and an inverse limit of their underlying groups. All profinite actions on sets are continuous for the discrete topology on the set. Cohomology of a topos has abelian coefficients. ``Torsion sheaf'' never means nonabelian cohomology in unbounded degrees. An expression with $\widehat{\mathbb Z}$ denotes continuous derived inverse-limit cohomology of its finite quotients, with the inverse-limit issue stated where used.

A geometric morphism $f:\mathcal E\to\mathcal F$ has inverse image $f^*:\mathcal F\to\mathcal E$. Pseudopullbacks and pseudopushouts retain invertible commutation data. A lax gluing category retains arbitrary specialization maps. Neither construction is interchangeable with its strict categorical counterpart.
The disjoint sum $\coprod_i\mathcal E_i$ of topoi has underlying category $\prod_i\mathcal E_i$; it is not their product in the two-category of topoi. Our diagrams use the former notation for disjoint strata.

\begin{longtable}{p{0.18\linewidth}p{0.25\linewidth}p{0.46\linewidth}}
\toprule
Our notation & Manuscript notation & Introduction and meaning\\
\midrule
$\beta G$ & $\operatorname{Ens}(G)$ & LM \S1, scans 6--11: discrete action topos; a continuous-action variant is expressly marked.\\
$\mathcal B(X,S)$ & $B_{X,U}$, $B_{U,X}$, $B_{X;S}$ & LM \S14, scans 123--127; \S16, scans 163--168: the mixed Artin gluing.\\
$G$, $H_s$, $Q_s$, $I_s$ & $\mathcal E$, $\mathcal E_i$, $\Gamma_i$, $\Pi_i$ & LM \S14, scans 122--123: global, local, residue and inertia groups.\\
$X^{*\infty}$ & $\mathcal V^*_{\Theta,X}$ & Cote 147, scan 150: punctured-neighborhood pro-object, with its precise realization specified.\\
\bottomrule
\end{longtable}
We cite LM by section, manuscript folio when marked, and scan. Scan numbers refer to the attached archive images, not to pages of this paper. Readings remain provisional; mathematical consistency is evidence, not a declaration of palaeographic certainty.
